\textbf{Пример 1}

Рассмотрим следующий вызов функции:

\texttt{construct\_roads([4, 4, 6, 4, 2], [4, 6, 4, 2, 4])}

Это означает, что в парке находятся $5$ фонтанов:
\begin{itemize}
\item фонтан $0$ расположен в точке $(4,4)$,
\item фонтан $1$ расположен в точке $(4,6)$,
\item фонтан $2$ расположен в точке $(6,4)$,
\item фонтан $3$ расположен в точке $(4,2)$,
\item фонтан $4$ расположен в точке $(2,4)$.
\end{itemize}

Можно проложить следующие $4$ дорожки, где каждая дорожка соединяет два фонтана, и разместить на них лавочки следующим образом:


\begin{center}
\renewcommand{\arraystretch}{1.5}
\begin{tabular}{|c|c|c|}
\hline
Номер дорожки & Фонтаны на концах & Расположение лавочки \\
\hline
$0$&$0,2$&$(5,5)$\\
\hline
$1$&$0,1$& $(3,5)$\\
\hline
$2$&$3,0$& $(5,3)$\\
\hline
$3$&$4,0$& $(3,3)$\\
\hline
\end{tabular}
\end{center}


Это решение соответствует следующей картинке:

\includegraphics[scale=0.7]{1.png}

Чтобы описать это решение, \texttt{construct\_roads} должна сделать следующий вызов:
\begin{itemize}
\item \texttt{build([0, 0, 3, 4], [2, 1, 0, 0], [5, 3, 5, 3], [5, 5, 3, 3])}
\end{itemize}

Затем необходимо вернуть значение $1$.

Обратите внимание, что в этом примере есть несколько возможных решений, удовлетворяющих ограничениям, любое из них будет принято.
Например, также можно описать решение следующим вызовом: \texttt{build([1, 2, 3, 4], [0, 0, 0, 0], [5, 5, 3, 3], [5, 3, 3, 5])}, а затем вернуть значение $1$.

\textbf{Пример 2}

Рассмотрим следующий вызов функции:

\texttt{construct\_roads([2, 4], [2, 6])}

Фонтан $0$ расположен в точке $(2,2)$, а фонтан $1$ расположен в точке $(4,6)$. Поскольку проложить дорожки в соответствии с описанными ограничениями нельзя, функция \texttt{construct\_roads} должна вернуть $0$, не вызывая перед этим \texttt{build}.
